% Einheit 4 – Ausklammern (bank/terme/e4.jsonl, zusammenbau v1.0)
\einheitenkopf[e4][Ausklammern]{Ausklammern}
\setcounter{aufgabe}{19}


% test: terme-e4-k2-s10-v2
\begin{kbaufgabe}{Kannst du das schon? Dann weiter zu „Terme aufstellen und berechnen“}
\begin{kbblock}
\kbanweisung{Klammere so weit wie möglich aus und mache die Probe.}
\kbteil{}{$5x^2 + 10xy - 15x$}{}
\item[]\kbraum{2}
\end{kbblock}
\end{kbaufgabe}


% erkennung: terme-e4-k1-s0-v1, terme-e4-k1-s0-v2, terme-e4-k1-s0-v3, terme-e4-k1-s0-v4
\begin{kbaufgabe}{Erkennen, welcher Faktor in jedem Glied steckt}
\begin{kbblock}
\kbteil{a)}{Was steckt in jedem Glied von $8x + 12$? Kreise den Faktor ein, der in beiden Gliedern steckt. Klammere nichts aus.}{}
\end{kbblock}
\begin{kbblock}
\kbteil{b)}{Was steckt in jedem Glied von $9y^2 + 2y$? Kreise den Faktor ein, der in beiden Gliedern steckt. Klammere nichts aus.}{}
\end{kbblock}
\begin{kbblock}
\kbteil{c)}{Was steckt in jedem Glied von $6a^2 + 10a$? Kreise ein, was in beiden Gliedern steckt: Zahl, Variable oder beides. Klammere nichts aus.}{}
\end{kbblock}
\begin{kbblock}
\kbteil{d)}{Was steckt in jedem Glied von $9x - 6 + 15x^2$? Kreise ein, was in allen drei Gliedern steckt. Klammere nichts aus.}{}
\end{kbblock}
\end{kbaufgabe}


% vorstufe: terme-e4-k2-s-1-v1, terme-e4-k2-s-1-v2, terme-e4-k2-s-1-v3, terme-e4-k2-s-1-v4
\begin{kbaufgabe}{Jedes Glied als Produkt mit einem Faktor schreiben}
\begin{kbblock}
\kbteil{a)}{Schreibe jedes Glied von $8x + 20$ als Produkt mit dem Faktor 4.}{}
\kbantwort{$4 \cdot$ \kbfeld $+ 4 \cdot$ \kbfeld}
\end{kbblock}
\begin{kbblock}
\kbteil{b)}{Schreibe jedes Glied von $6a + 14$ als Produkt mit dem Faktor 2.}{}
\kbantwort{$2 \cdot$ \kbfeld $+ 2 \cdot$ \kbfeld}
\end{kbblock}
\begin{kbblock}
\kbteil{c)}{Schreibe jedes Glied von $10y + 15$ als Produkt mit dem Faktor 5.}{}
\kbantwort{$5 \cdot$ \kbfeld $+ 5 \cdot$ \kbfeld}
\end{kbblock}
\begin{kbblock}
\kbteil{d)}{Schreibe jedes Glied von $9b + 6$ als Produkt mit dem Faktor 3.}{}
\kbantwort{$3 \cdot$ \kbfeld $+ 3 \cdot$ \kbfeld}
\end{kbblock}
\end{kbaufgabe}


% vorstufe: terme-e4-k2-s0-v1, terme-e4-k2-s0-v2, terme-e4-k2-s0-v3, terme-e4-k2-s0-v4
\begin{kbaufgabe}{Einen vorgegebenen Faktor ausklammern}
\begin{kbblock}
\lbpaar{\kbteil{a)}{Klammere den Faktor 3 aus: $9x + 21$.}{}
\kbantwort{$3 \cdot ($ \kbfeld $+$ \kbfeld $)$}
}{\kbteil{b)}{Klammere den Faktor 2 aus: $4y + 10$.}{}
\kbantwort{$2 \cdot ($ \kbfeld $+$ \kbfeld $)$}}
\end{kbblock}
\begin{kbblock}
\lbpaar{\kbteil{c)}{Klammere den Faktor 5 aus: $10a + 35$.}{}
\kbantwort{$5 \cdot ($ \kbfeld $+$ \kbfeld $)$}
}{\kbteil{d)}{Klammere den Faktor 7 aus: $14b + 21$.}{}
\kbantwort{$7 \cdot ($ \kbfeld $+$ \kbfeld $)$}}
\end{kbblock}
\end{kbaufgabe}


% leiter: terme-e4-k2-s1-v1, terme-e4-k2-s1-v2, terme-e4-k2-s1-v3, terme-e4-k2-s1-v4, terme-e4-k2-s2-v1, terme-e4-k2-s3-v1, terme-e4-k2-s4-v1, terme-e4-k2-s5-v1, terme-e4-k2-s6-v1, terme-e4-k2-s9-v1
\begin{kbaufgabe}{Ausklammern}
\begin{kbblock}
\kbanweisung{Klammere den größten gemeinsamen Faktor aus.}
\lbpaar{\kbteil{a)}{$5x + 10 =$ \lbfeld}{}
}{\kbteil{b)}{$5x + 25 =$ \lbfeld}{}}
\end{kbblock}
\begin{kbblock}
\lbpaar{\kbteil{c)}{$5x + 30 =$ \lbfeld}{}
}{\kbteil{d)}{$5x + 35 =$ \lbfeld}{}}
\end{kbblock}
\begin{kbblock}
\lbpaar{\kbteil{e)}{$x^2 + 5x =$ \lbfeld}{}
}{\kbteil{f)}{$6x^2 + 15x =$ \lbfeld}{}}
\end{kbblock}
\begin{kbblock}
\lbpaar{\kbteil{g)}{$6x + 6 =$ \lbfeld}{}
}{\kbteil{h)}{$3x^2 + 6x + 15 =$ \lbfeld}{}}
\end{kbblock}
\begin{kbblock}
\kbanweisung{Klammere aus und mache die Probe.}
\kbteil{i)}{$6x + 30$}{}
\item[]\kbraum{2}
\end{kbblock}
\begin{kbblock}
\kbteil{j)}{Nora soll im Term $5x + 20$ den größten gemeinsamen Faktor ausklammern. Sie rechnet so: \rechnung{5x + 20 &= 5 \cdot (x + 20)} Finde den Fehler und rechne richtig.}{}
\item[]\kbraum{2}
\end{kbblock}
\end{kbaufgabe}


% pruefung: terme-e4-k2-s10-v1
\begin{kbaufgabe}{Auch Terme mit zwei Variablen so weit wie möglich ausklammern}
\begin{kbblock}
\kbanweisung{Klammere so weit wie möglich aus und mache die Probe.}
\kbteil{}{$3a^2 + 6ab + 9a$}{}
\item[]\kbraum{2}
\end{kbblock}
\end{kbaufgabe}


% pflicht: terme-e4-k3-s1-v1, terme-e4-k3-s2-v1, terme-e4-k3-s3-v1, terme-e4-k3-s4-v1
\begin{kbaufgabe}{Verstanden?}
\begin{kbblock}
\kbteil{a)}{Finn soll im Term $20x - 8$ den größten gemeinsamen Faktor ausklammern. Er rechnet so: \rechnung{20x - 8 &= 4 \cdot (5x + 2)} Finde den Fehler und rechne richtig.}{}
\item[]\kbraum{2}
\end{kbblock}
\begin{kbblock}
\kbteil{b)}{Begründe, warum die Terme $6a + 42$ und $6 \cdot (a + 7)$ für jede Zahl $a$ denselben Wert haben. Prüfe es auch durch Einsetzen von $a = 2$.}{}
\item[]\kbraum{3}
\end{kbblock}
\begin{kbblock}
\kbteil{c)}{Der Term $6 \cdot (b + 3)$ gibt den Flächeninhalt einer Figur aus zwei Rechtecken in cm² an. Beschreibe die Figur in Worten, ohne sie zu zeichnen.}{}
\item[]\kbraum{2}
\end{kbblock}
\begin{kbblock}
\kbteil{d)}{Ein Sportverein kauft 15 Trikots zu je 28 € und 15 Hosen zu je 17 €. Schreibe die Kosten als Term mit Klammer. Der Verein hat 700 €. Reicht das Geld?}{}
\item[]\kbraum{2}
\end{kbblock}
\end{kbaufgabe}
